\documentclass{article}

\usepackage{amsmath,amssymb}
\usepackage{xurl}
\usepackage[hidelinks]{hyperref}
\setlength{\emergencystretch}{2em}

% Shared commands for notes in docs/paper-gaps/.

\DeclareUnicodeCharacter{2080}{\ifmmode{}_{0}\else\textsubscript{0}\fi}
\DeclareUnicodeCharacter{2081}{\ifmmode{}_{1}\else\textsubscript{1}\fi}
\DeclareUnicodeCharacter{2082}{\ifmmode{}_{2}\else\textsubscript{2}\fi}
\DeclareUnicodeCharacter{2099}{\ifmmode{}_{n}\else\textsubscript{n}\fi}

\newcommand{\Lean}{\textsc{Lean}}
\ExplSyntaxOn
\NewDocumentCommand{\leanid}{m}{
  \ifmmode
    \textsf{\detokenize{#1}}
  \else
    \group_begin:
    \urlstyle{sf}
    \tl_set:Nn \l_tmpa_tl {#1}
    \tl_replace_all:Nnn \l_tmpa_tl {\_} {_}
    \exp_args:NV \path \l_tmpa_tl
    \group_end:
  \fi
}
\ExplSyntaxOff
\providecommand{\tr}{\operatorname{tr}}
\newcommand{\tnleanIssueBase}{https://github.com/LionSR/TNLean/issues/}
\newcommand{\tnleanPrBase}{https://github.com/LionSR/TNLean/pull/}
\newcommand{\tnleanDocBase}{https://sirui-lu.com/TNLean/docs/}
\newcommand{\ghissue}[1]{\href{\tnleanIssueBase#1}{GitHub issue \##1}}
\newcommand{\ghpr}[1]{\href{\tnleanPrBase#1}{PR \##1}}
\newcommand{\doclink}[1]{\href{\tnleanDocBase#1.html}{\path{#1}}}

% Dirac notation and the identity operator, as in blueprint/src/macros/common.tex.
% \providecommand keeps note-local definitions authoritative.
\usepackage{dsfont}
\providecommand{\ket}[1]{|#1\rangle}
\providecommand{\bra}[1]{\langle #1|}
\providecommand{\braket}[2]{\langle #1 | #2 \rangle}
\providecommand{\ketbra}[2]{|#1\rangle\!\langle #2|}
\providecommand{\Id}{\mathds{1}}
\providecommand{\one}{\mathds{1}}
\providecommand{\C}{\mathbb{C}}
\providecommand{\MN}[1]{M_{#1}(\C)}
\providecommand{\MD}{M_D(\C)}

% External referencing for paper sources.  The build helper writes sanitized
% auxiliary files under build/paper-aux/ and excludes duplicated source labels.
\usepackage{xr}

% Import a generated paper auxiliary file with a caller-supplied label prefix.
% The candidate paths support builds from docs/paper-gaps/, the repository root,
% and build/paper-aux/ itself.
\newcommand{\importpaperaux}[2]{%
  \IfFileExists{../../build/paper-aux/#2.aux}{%
    \externaldocument[#1]{../../build/paper-aux/#2}%
  }{%
    \IfFileExists{build/paper-aux/#2.aux}{%
      \externaldocument[#1]{build/paper-aux/#2}%
    }{%
      \IfFileExists{#2.aux}{%
        \externaldocument[#1]{#2}%
      }{}%
    }%
  }%
}

% General external reference macros
\newcommand{\paperref}[2]{\ref{#1-#2}}
\newcommand{\papereqref}[2]{(\ref{#1-#2})}

% CPSV16 source (1606.00608 / MPDO-22-12-17-2.tex) external document and shortcuts
\importpaperaux{cpsv16-}{1606.00608/MPDO-22-12-17-2-tnlean}
\newcommand{\cpsvref}[1]{\paperref{cpsv16}{#1}}
\newcommand{\cpsveqref}[1]{\papereqref{cpsv16}{#1}}

% Verdict marker: \gapnote{<kind>}{<status>}. Kind is the policy
% classification (clarification, local-correction, scope-restriction,
% unfaithful, false-source, open-gap); status is open, wip (elimination
% underway), resolved, or historical. Machine-read by the site index;
% prints a verdict line.
\newcommand{\gapnote}[2]{\par\noindent\textbf{Verdict:} #1 (#2).\par}


\title{The Square-Dimension Restriction in Wolf's Generic Normal Form}
\author{}
\date{2026-06-30}

\begin{document}
\maketitle

\section{The assertion}

Wolf's Proposition~2.9 is a normal-form statement for completely positive maps
with full Kraus rank.\footnote{\cite[Proposition~2.9]{Wolf2012Quantum};
local source \path{Notes/WolfNoteTexSource/ch02_representations.tex},
lines 923--929.}  In the source it reads: let
\begin{align}
  T : \mathcal{M}_{d_1} \to \mathcal{M}_{d_2}
\end{align}
be a completely positive map with full Kraus rank.  Then there exist invertible
completely positive maps \(\Phi_i \in \mathcal B(\mathcal{M}_{d_i})\) of Kraus
rank one such that \(T' := \Phi_2 T \Phi_1\) is doubly-stochastic, i.e.\ both
\(T'(\one)\) and \(T'^{*}(\one)\) are proportional to the identity matrix.

The two filterings act on \emph{different} matrix algebras: \(\Phi_1\) on the
input algebra \(\mathcal{M}_{d_1}\) and \(\Phi_2\) on the output algebra
\(\mathcal{M}_{d_2}\), with no constraint relating \(d_1\) and \(d_2\).

\section{The point at issue}

The formal development states the proposition only in the equal-dimension case
\(d_1 = d_2 = D\):
\begin{align}
  T : \MN{D} \to \MN{D},
\end{align}
with both filterings acting on the same algebra \(\MN{D}\).\footnote{Formal
declaration \leanid{Wolf.exists\_normal\_form\_generic} in
\doclink{TNLean/Channel/LorentzNormalForm}.}  The Choi matrix
\(\tau = (T \otimes \mathrm{id})(\ketbra{\Omega}{\Omega})\) is then a square
matrix on \(\C^{D} \otimes \C^{D}\), and the minimisation runs over
\(\mathrm{SL}(D,\C) \times \mathrm{SL}(D,\C)\) filterings of a single dimension.
The compactness lemma \leanid{Wolf.infimum\_is\_attained} and the
trace--determinant AM--GM equality
\leanid{Matrix.PosSemidef.pow\_card\_mul\_det\_re\_eq\_trace\_re\_pow\_iff}
that drive the proof are likewise written for a square Choi matrix on
\(\C^{D} \otimes \C^{D}\).

For \(D=2\), the square theorem supplies the full-Kraus-rank minimisation result
used in the diagonal case of the qubit Lorentz normal form.  The existence theorem
covering the diagonal, non-diagonal, and singular cases is not yet formalized; in
particular, it is not a specialization of a current declaration.  Its additional
scalar-normalization and Lorentz-orbit gap is described in
\path{docs/paper-gaps/wolf_prop2_11_lorentz_scalar_filtering_gap.tex}.  The general
rectangular statement of Proposition~2.9 is also not yet formalized; only its
equal-dimension instance is.

\section{The restricted statement}

The formal theorem proves the square instance:
\begin{align}
  T : \MN{D} \to \MN{D}, \qquad
  \tau = \mathrm{choi}(T) \succ 0
  \ \Longrightarrow\
  \exists\, \Phi_1, \Phi_2,\
  \Phi_2 \circ T \circ \Phi_1\ \text{doubly-stochastic},
\end{align}
with both \(\Phi_1, \Phi_2\) SL-filterings on \(\MN{D}\).  This is a sub-case of
the source proposition, obtained by setting \(d_1 = d_2\); it is mathematically
correct as stated, just narrower than Wolf's rectangular formulation.

\section{Elimination plan}

A faithful formalization of the source proposition requires generalising the
two supporting lemmas to independent input and output dimensions:
\begin{itemize}
  \item \leanid{Wolf.infimum\_is\_attained} should range over a Choi matrix on
    \(\C^{d_2} \otimes \C^{d_1}\) with filterings in
    \(\mathrm{SL}(d_2,\C) \times \mathrm{SL}(d_1,\C)\);
  \item \leanid{Wolf.exists\_normal\_form\_generic} should then be restated for
    \(T : \mathcal{M}_{d_1} \to \mathcal{M}_{d_2}\) with separate filterings on
    each factor.
\end{itemize}
Until that work is done, the square version is kept as an explicitly-labelled
specialisation, and the source-faithful rectangular proposition remains
unformalised.  The blueprint entry \texttt{thm:exists\_normal\_form\_generic}
and the Lean docstring carry the matching scope marker.

\bibliographystyle{alpha}
\bibliography{references}

\end{document}
